\documentclass[11pt]{article}
\usepackage[a4paper,landscape,margin=15mm,top=18mm,bottom=16mm]{geometry}
\usepackage{fontspec}
\setmainfont{Arial}
\usepackage{graphicx}
\usepackage{xcolor}
\usepackage{fancyhdr}
\usepackage{sectsty}
\usepackage{hyperref}
\definecolor{eltongold}{RGB}{169,138,87}
\definecolor{eltongold2}{RGB}{208,174,120}
\definecolor{eltondark}{RGB}{15,15,15}
\definecolor{eltonred}{RGB}{255,51,51}
\hypersetup{colorlinks=true,linkcolor=eltongold!80!black}
\renewcommand{\contentsname}{Contents}
\sectionfont{\color{eltongold!75!black}}
\subsectionfont{\color{eltondark}}
\pagestyle{fancy}\fancyhf{}
\fancyhead[L]{\small\color{eltongold!75!black}\textbf{ELTON} \color{gray}PDM25 V2 — Wiring manual}
\fancyhead[R]{\small\color{gray}2026-06-17}
\fancyfoot[C]{\small\color{gray}\thepage}
\renewcommand{\headrulewidth}{1pt}
\renewcommand{\headrule}{\hbox to\headwidth{\color{eltongold}\leaders\hrule height 1pt\hfill}}
\begin{document}
\thispagestyle{empty}
\pagecolor{eltondark}
\vspace*{1.4cm}\begin{center}
\includegraphics[width=0.52\textwidth]{C:/Users/joel/Documents/Elton LT Wire diagrams/hwpdm-project/docs/wireviz_v8.0/_assets/elton-logo.png}\\[14pt]
{\fontsize{42}{48}\selectfont\bfseries\color{eltongold2} ELTON}\\[8pt]
{\color{eltongold}\rule{60mm}{1.2pt}}\\[10pt]
{\LARGE\color{white} Electrical system \& wiring}\\[4pt]
{\large\color{white!80} Wiring diagrams — service manual}\\[12pt]
{\large\color{eltongold} PDM25 V2 \quad·\quad VW LT31 1976 \quad·\quad JSN 398}\\[14pt]
{\normalsize\color{white!75} Prepared by Joel Kvarnsmyr}\\[3pt]
{\small\color{white!55} 2026-06-17}
\end{center}\vfill
\clearpage
\pagecolor{white}
\thispagestyle{empty}
\tableofcontents\clearpage
\section{Introduction}
This is the physical wiring and service manual for the conversion of a VW LT31 (1976, JSN~398) to a Hardwire PDM25 V2 + Raspberry Pi dashboard. Each diagram shows the PDM connector on the left, the cable in the middle (cross-section + actual color) and the consumer on the right.\\[6pt]\textbf{Conventions:} The connectors' title bar carries the housing's physical color (A grey, B black, C green, D brown). Cross-sections are given in mm\textsuperscript{2}. Ground is shown as a GND pin on the consumer rather than as a drawn cable. Wire colors reflect the wiring as physically inspected in the vehicle.
\clearpage
\section{Connectors}
\subsection{Connector A — Inputs (GREY)}
{\small All input signals: levers, switches and senders. Switch-to-ground unless otherwise noted. Grey housing. Pin number on the left, cable in the middle (cross-section + actual color), consumer on the right.}
\\[4pt]
\begin{center}
\includegraphics[width=\linewidth,height=0.74\textheight,keepaspectratio]{C:/Users/joel/Documents/Elton LT Wire diagrams/hwpdm-project/docs/wireviz_v8.0/1_kontakter/kontakt_A_inputs.png}
\end{center}
{\scriptsize\color{gray} Diagram: wireviz\_v8.0/1\_kontakter/kontakt\_A\_inputs}
\clearpage
\subsection{Connector B — Outputs (BLACK)}
{\small Outputs: low beam, left turn signal, horn 1, wiper low, blower, washer pump, starter motor, reversing light, USB.}
\\[4pt]
\begin{center}
\includegraphics[width=\linewidth,height=0.74\textheight,keepaspectratio]{C:/Users/joel/Documents/Elton LT Wire diagrams/hwpdm-project/docs/wireviz_v8.0/1_kontakter/kontakt_B_outputs.png}
\end{center}
{\scriptsize\color{gray} Diagram: wireviz\_v8.0/1\_kontakter/kontakt\_B\_outputs}
\clearpage
\subsection{Connector C — Power / CAN / sensors (GREEN)}
{\small Control feed (C3), chassis ground, CAN to the Pi, ignition/cranking, ignition coil, alternator field excitation, signal busbar and the 5V reference distribution point.}
\\[4pt]
\begin{center}
\includegraphics[width=\linewidth,height=0.74\textheight,keepaspectratio]{C:/Users/joel/Documents/Elton LT Wire diagrams/hwpdm-project/docs/wireviz_v8.0/1_kontakter/kontakt_C_kraft_can.png}
\end{center}
{\scriptsize\color{gray} Diagram: wireviz\_v8.0/1\_kontakter/kontakt\_C\_kraft\_can}
\clearpage
\subsection{Connector D — Outputs (BROWN)}
{\small Outputs: high beam, parking light, brake light, right turn signal, fuel pump (2 parallel to the same pin), wiper high, horn 2.}
\\[4pt]
\begin{center}
\includegraphics[width=\linewidth,height=0.74\textheight,keepaspectratio]{C:/Users/joel/Documents/Elton LT Wire diagrams/hwpdm-project/docs/wireviz_v8.0/1_kontakter/kontakt_D_outputs.png}
\end{center}
{\scriptsize\color{gray} Diagram: wireviz\_v8.0/1\_kontakter/kontakt\_D\_outputs}
\clearpage
\section{Power supply}
\subsection{+12V power busbar (always-on loads)}
{\small The busbar (battery-fed via a 10A ceramic fuse) supplies the always-on loads: radio and Raspberry Pi. The PDM control feed C3 has its own graph (next page).}
\\[4pt]
\begin{center}
\includegraphics[width=\linewidth,height=0.74\textheight,keepaspectratio]{C:/Users/joel/Documents/Elton LT Wire diagrams/hwpdm-project/docs/wireviz_v8.0/2_kraft/kraft_busbar.png}
\end{center}
{\scriptsize\color{gray} Diagram: wireviz\_v8.0/2\_kraft/kraft\_busbar}
\clearpage
\subsection{PDM control feed C3 — ignition + hazard wake}
{\small Two paths to C3 from the busbar: normally via the ignition switch (RUN), and as a wake via the hazard switch through a blocking diode. Vertical layout.}
\\[4pt]
\begin{center}
\includegraphics[width=\linewidth,height=0.74\textheight,keepaspectratio]{C:/Users/joel/Documents/Elton LT Wire diagrams/hwpdm-project/docs/wireviz_v8.0/2_kraft/kraft_c3.png}
\end{center}
{\scriptsize\color{gray} Diagram: wireviz\_v8.0/2\_kraft/kraft\_c3}
\clearpage
\subsection{5V reference — analog inputs}
{\small C10 (5V REF, max 100 mA) feeds, via a distribution point, five SEPARATE pull-ups to the analog inputs (wiper/blower/turn mux plus coolant and fuel). Total load \textasciitilde{}37 mA. Each input has its own pull-up — never shared.}
\\[4pt]
\begin{center}
\includegraphics[width=\linewidth,height=0.74\textheight,keepaspectratio]{C:/Users/joel/Documents/Elton LT Wire diagrams/hwpdm-project/docs/wireviz_v8.0/2_kraft/kraft_5v.png}
\end{center}
{\scriptsize\color{gray} Diagram: wireviz\_v8.0/2\_kraft/kraft\_5v}
\clearpage
\section{Circuits}
\subsection{Mux — wiper + washer (I9): switch + resistors}
{\small The resistor coding for the wiper+washer mux. Washer 10k, low 4.7k, high 22k on the switch outputs → chassis ground; each position gives a unique node voltage. Feed/sensing (5.6k pull-up → 5V REF) is shown in the 5V-reference diagram.}
\\[4pt]
\begin{center}
\includegraphics[width=\linewidth,height=0.74\textheight,keepaspectratio]{C:/Users/joel/Documents/Elton LT Wire diagrams/hwpdm-project/docs/wireviz_v8.0/3_kretsar/mux_torkare_brytare.png}
\end{center}
{\scriptsize\color{gray} Diagram: wireviz\_v8.0/3\_kretsar/mux\_torkare\_brytare}
\clearpage
\subsection{Mux — blower low/high (I1): knob + resistors}
{\small The blower knob's resistor coding. Low 10k, high 2.2k → chassis ground. Feed via pull-up to 5V REF: see the 5V-reference diagram. Two speeds come from O11's PWM — no power resistor.}
\\[4pt]
\begin{center}
\includegraphics[width=\linewidth,height=0.74\textheight,keepaspectratio]{C:/Users/joel/Documents/Elton LT Wire diagrams/hwpdm-project/docs/wireviz_v8.0/3_kretsar/mux_flakt_brytare.png}
\end{center}
{\scriptsize\color{gray} Diagram: wireviz\_v8.0/3\_kretsar/mux\_flakt\_brytare}
\clearpage
\subsection{Coolant temperature — NTC (I10)}
{\small NTC sender as a voltage divider against an external pull-up to 5V REF (\textasciitilde{}470-680Ω). High temperature = low resistance = low voltage. Ground via the sender thread in the engine block.}
\\[4pt]
\begin{center}
\includegraphics[width=\linewidth,height=0.74\textheight,keepaspectratio]{C:/Users/joel/Documents/Elton LT Wire diagrams/hwpdm-project/docs/wireviz_v8.0/3_kretsar/vattentemp.png}
\end{center}
{\scriptsize\color{gray} Diagram: wireviz\_v8.0/3\_kretsar/vattentemp}
\clearpage
\subsection{Oil-pressure switch (I14)}
{\small Simple switch-to-ground (F1). Internal pull-up to VBat. Low pressure / engine off = 0V (warning), engine running + pressure OK = high.}
\\[4pt]
\begin{center}
\includegraphics[width=\linewidth,height=0.74\textheight,keepaspectratio]{C:/Users/joel/Documents/Elton LT Wire diagrams/hwpdm-project/docs/wireviz_v8.0/3_kretsar/oljetryck.png}
\end{center}
{\scriptsize\color{gray} Diagram: wireviz\_v8.0/3\_kretsar/oljetryck}
\clearpage
\subsection{Handbrake warning (I16)}
{\small Mechanical switch-to-ground (F9), grounds when the handbrake is applied. Internal pull-up, a copy of the brake-light input.}
\\[4pt]
\begin{center}
\includegraphics[width=\linewidth,height=0.74\textheight,keepaspectratio]{C:/Users/joel/Documents/Elton LT Wire diagrams/hwpdm-project/docs/wireviz_v8.0/3_kretsar/handbroms.png}
\end{center}
{\scriptsize\color{gray} Diagram: wireviz\_v8.0/3\_kretsar/handbroms}
\clearpage
\subsection{Ignition coil (1/2): feed O3 → N6 → coil}
{\small Original VW, but fed by O3 instead of the ignition switch. O3 → N6 ballast resistor → coil Term. 15.}
\\[4pt]
\begin{center}
\includegraphics[width=\linewidth,height=0.74\textheight,keepaspectratio]{C:/Users/joel/Documents/Elton LT Wire diagrams/hwpdm-project/docs/wireviz_v8.0/3_kretsar/motor_tandspole_1_matning.png}
\end{center}
{\scriptsize\color{gray} Diagram: wireviz\_v8.0/3\_kretsar/motor\_tandspole\_1\_matning}
\clearpage
\subsection{Ignition coil (2/2): coil → distributor}
{\small Coil Term. 1 → distributor breaker points → ground. The points break the primary circuit → HT is induced (Term. 4 → distributor cap → spark plug, not drawn).}
\\[4pt]
\begin{center}
\includegraphics[width=\linewidth,height=0.74\textheight,keepaspectratio]{C:/Users/joel/Documents/Elton LT Wire diagrams/hwpdm-project/docs/wireviz_v8.0/3_kretsar/motor_tandspole_2_fordelare.png}
\end{center}
{\scriptsize\color{gray} Diagram: wireviz\_v8.0/3\_kretsar/motor\_tandspole\_2\_fordelare}
\clearpage
\subsection{Alternator (1/2): field excitation}
{\small O2 → 82Ω power resistor → blocking diode → alternator D+. The resistor supplies field current; the diode stops D+ (\textasciitilde{}14V) back-feeding the PDM.}
\\[4pt]
\begin{center}
\includegraphics[width=\linewidth,height=0.74\textheight,keepaspectratio]{C:/Users/joel/Documents/Elton LT Wire diagrams/hwpdm-project/docs/wireviz_v8.0/3_kretsar/generator_1_excitering.png}
\end{center}
{\scriptsize\color{gray} Diagram: wireviz\_v8.0/3\_kretsar/generator\_1\_excitering}
\clearpage
\subsection{Alternator (2/2): charge status}
{\small Charge status is tapped AFTER the diode to a spare PDM input (charging \textasciitilde{}14V = high, not charging \textasciitilde{}1V = low) → CAN → charge warning lamp in the dashboard.}
\\[4pt]
\begin{center}
\includegraphics[width=\linewidth,height=0.74\textheight,keepaspectratio]{C:/Users/joel/Documents/Elton LT Wire diagrams/hwpdm-project/docs/wireviz_v8.0/3_kretsar/generator_2_laddstatus.png}
\end{center}
{\scriptsize\color{gray} Diagram: wireviz\_v8.0/3\_kretsar/generator\_2\_laddstatus}
\clearpage
\section{Dashboard}
\subsection{Dashboard (1/2): Pi → GX16 rainbow ribbon}
{\small Pi 5 → GX16 10-pin ribbon in fixed color order. Canonical GPIO source: DRIVER-DASH.md §7.}
\\[4pt]
\begin{center}
\includegraphics[width=\linewidth,height=0.74\textheight,keepaspectratio]{C:/Users/joel/Documents/Elton LT Wire diagrams/hwpdm-project/docs/wireviz_v8.0/4_dash/pi_1_ribbon.png}
\end{center}
{\scriptsize\color{gray} Diagram: wireviz\_v8.0/4\_dash/pi\_1\_ribbon}
\clearpage
\subsection{Dashboard (2/2): GX16 → 3× ST7789 (SPI0)}
{\small Pins 1-7 (GND/3V3/SCLK/MOSI/RST/DC/BLK) are daisy-chained to all three displays (shown as a shared bus); only CS is unique: orange→S1, red→S2, brown→S3.}
\\[4pt]
\begin{center}
\includegraphics[width=\linewidth,height=0.74\textheight,keepaspectratio]{C:/Users/joel/Documents/Elton LT Wire diagrams/hwpdm-project/docs/wireviz_v8.0/4_dash/pi_2_displays.png}
\end{center}
{\scriptsize\color{gray} Diagram: wireviz\_v8.0/4\_dash/pi\_2\_displays}
\clearpage
\section{Appendix A — Component reference}
{\small Components on the converted vehicle, by their VW designation as used in the diagrams. From the VW LT workshop manual (Fig.\ 13.77, 4-cyl).}\\[6pt]
\begin{center}\small
\begin{tabular}{|l|p{78mm}|l|p{78mm}|}
\hline
\textbf{ID} & \textbf{Description} & \textbf{ID} & \textbf{Description} \\ \hline
A & Battery (12 V) & H1 & Horn \\ \hline
B & Starter motor & H2 & Horn 2 (additional) \\ \hline
C & Alternator & L1 & Headlight, left (twin-filament) \\ \hline
D & Ignition / start switch & L2 & Headlight, right (twin-filament) \\ \hline
E2 & Turn-signal stalk & M5 & Turn signal, front left \\ \hline
E4 & High-beam / headlight-flasher stalk & M6 & Turn signal, rear left \\ \hline
E9 & Blower (fresh-air fan) switch & M7 & Turn signal, front right \\ \hline
E22 & Wiper / washer stalk & M8 & Turn signal, rear right \\ \hline
F & Brake-light switch & M9 & Brake light, left \\ \hline
F1 & Oil-pressure switch & M10 & Brake light, right \\ \hline
F4 & Reversing-light switch & M16 & Reversing light, left \\ \hline
F9 & Handbrake warning switch & M17 & Reversing light, right \\ \hline
G & Fuel-gauge sender & N & Ignition coil \\ \hline
G2 & Coolant-temperature sender & N6 & Series (ballast) resistor \\ \hline
G6 & Electric fuel pump & O & Distributor (breaker points) \\ \hline
\end{tabular}
\end{center}
\clearpage
\section{Appendix B — Terminal designations (DIN 72552)}
{\small Standard terminal numbers used throughout the diagrams.}\\[6pt]
\begin{center}\small
\begin{tabular}{|l|p{78mm}|l|p{78mm}|}
\hline
\textbf{ID} & \textbf{Description} & \textbf{ID} & \textbf{Description} \\ \hline
1 & Ignition coil primary, to distributor & 53 & Wiper, normal speed \\ \hline
4 & Ignition coil HT output (to distributor cap) & 53a & Wiper, fast speed \\ \hline
15 & Switched ignition + (live with ignition on) & 53b & Wiper park / return \\ \hline
30 & Battery + (permanent, unswitched) & 56a & High beam \\ \hline
31 & Ground (chassis earth return) & 56b & Low beam \\ \hline
50 & Starter control (solenoid, cranking only) & 58 & Parking / side lights \\ \hline
\end{tabular}
\end{center}
\vspace{6pt}\par{\footnotesize\color{gray} Terminal numbers follow the DIN 72552 standard and are common to all VW / Bosch wiring.}
\clearpage
\section{Appendix C — Earth points \& grounding}
{\small In the diagrams, ground is shown as a GND pin on each consumer rather than as a drawn cable. The physical chassis earth points below are where those grounds terminate. \textbf{All ground wires are brown.}}\\[6pt]
\begin{center}\small
\begin{tabular}{|l|p{70mm}|l|p{70mm}|}
\hline
\textbf{ID} & \textbf{Description} & \textbf{ID} & \textbf{Description} \\ \hline
1 & Battery earth strap to body & 14 & Roof cross-member, passenger side \\ \hline
3 & Engine-to-body earth strap & 15 & Rear, on longitudinal member \\ \hline
10 & Instrument-panel insert & 17 & At steering gear \\ \hline
11 & Behind instrument panel & 18 & Longitudinal member, rear left \\ \hline
12 & Under instrument panel (former fuse-box area) & 19 & Longitudinal member, rear right \\ \hline
13 & At rear door & 20 & In engine bay \\ \hline
\end{tabular}
\end{center}
\vspace{4pt}\par{\footnotesize\color{gray} Earth-point locations are per the VW LT factory manual for the donor vehicle and should be confirmed on the actual car. Common gauges: 0.5 mm\textsuperscript{2} (lamps, senders), 1.0 mm\textsuperscript{2} (motors, brake lights), 4.0 mm\textsuperscript{2} and heavier (battery, engine and gearbox straps).}
\clearpage
\end{document}